\scanpage{10}
% L01–03
\uline{Lemme} Prenons $k=\mathbb F_p[X_1,\text{--},X_n]$, $\xi_i=X_i$, donc les \deleted{\unsure{différents|systèmes}} $\Delta_i\in\mathbb F_p[X_1,\text{--},X_n]$ ($0\leq i\leq n$) \deleted{\unsure{donc}}, et on \deleted{Alors} a
% L04
\[ \Delta_n X_i^{p^n}+\Delta_{n-1}X_i^{p^{n-1}}\text{---}+\Delta_1X_i^p+\Delta_0X_i=0 \]
% L05–06
($1\leq i\leq n$). Les $\Delta_i$ ($0\leq i\leq n-1$) sont divisibles par $\Delta_n$, soit
% L07
\[ \text{et}\qquad \Delta_i=A_i\Delta_n \qquad 0\leq i\leq n-1 \]
% L08
\[ \text{et On a donc}\qquad \Delta_0=\pm\Delta_n^p(-1)^n\quad\text{au signe près} \]

% L09
\uline{Cor 1} \unsure{Ces|Cor}
% L09–10
\[ X_i^{p^n}+A_{n-1}\deleted{\unsure{X}}X_i^{p^{n-1}}\text{---}+A_1X_i^p\unsure{-|+} A_0X_i=0 \qquad (1\leq i\leq n) \]
% L11–12
\deleted{\unsure{et}} et les $A_i\in k[X_1,\text{--},X_n]$ ($0\leq i\leq n-1$) sont caractérisés par cette relation. \deleted{\unsure{Enfin|En effet}} De plus $A_0=\pm\Delta_n^{p-1}(-1)^n$ \deleted{\unsure{(au signe}}.

% L11–15
% margin note written sideways
\margin{\uline{Cor. 2} les $A_i$ sont \deleted{\unsure{stables}} \inserted{invariants} par l'action de $\operatorname{Gl}(n,\mathbb F_p)=\Gamma_n$ (a fortiori \deleted{et} \uline{symétriques}); $A_i$ est \uline{homogène} de degré $p^n-p^i$.}
% L13–15
\uline{Cor. 2} Revenons {:} une \deleted{$k$ qqcle. sur} \inserted{extn $k$ de} $\mathbb F_p$. Pour que l'on ait $\Delta_n(\xi)\neq0$ i.e. $\xi_1\text{--}\xi_n$ lin. indép sur $\mathbb F_p$, il f et s que $\Delta_0(\xi)\neq0$.

% L16
Considérons
% L16
\diagram{\begin{tikzcd}[ampersand replacement=\&]
\mathbb E^n \arrow[r,"A_*"] \& \mathbb E^n
\end{tikzcd}}
% L17
défini par $A_0\text{--}A_{n-1}$. Il se factorise en
% L18
\deleted{\unsure{Prop} \unsure{Cor 3} \unsure{Le $A_*$ induit un iso}}
% L19
\[ \deleted{U_n\,(=\mathbb E^n_{\Delta_n})=A_*^{-1}(\mathbb E^*\times\mathbb E^{n-1})} \]
% L20
\deleted{et $A_*$ induit un iso}
% L21–22
\diagram{\begin{tikzcd}[ampersand replacement=\&]
\mathbb E^n \arrow[rr,no head,"A_*"] \arrow[dr,no head,"\mathrm{can}" description] \& \& \mathbb E^n\\
\& \unsure{\mathbb Q|Q}_n \arrow[ur,"\psi_n"'] \&
\end{tikzcd}}
% L23
\deleted{et induit} Ceci dit.
% archive sheet number: 9
